The three things had been these.

\textit{Don’t let anyone move me.} That had been the first, and it had been the easiest, because there was nobody in the Tower of the Hand to move him — the whole castle had emptied downward hours ago, into the streets and the works and the hospital in the Great Hall, and the topmost room of the Tower of the Hand was the single quietest place in King’s Landing that night, which Sam had come to understand was very probably why Bran had chosen it.

\textit{Write down anything I say. Exactly as I say it, even if it’s nonsense. Especially if it’s nonsense.} That had been the second, and Sam had said yes to it before Bran had finished the sentence, and had gone and got the wax tablet and the stylus from the small library off the Tower rather than paper and ink, because paper and ink required a hand that was not shaking and wax did not.

The third one Bran had not finished.

“If Meera Reed ever comes here,” he had said, and stopped.

Sam had waited. He was good at waiting; it was very nearly the only thing he had ever been unambiguously good at. He had stood there with the tablet under his arm and waited while Bran Stark, who had not hesitated over a sentence in two years, sat in a chair at the top of a tower with his hands in his lap and did not find the end of one.

“Bran?”

“I don’t have it,” Bran had said. He had sounded, for perhaps a second and a half, almost exactly like a boy. “There isn’t a thing to tell her. I’ve been looking for one for four hundred and six days and there isn’t one, and I’d rather you knew that than went and made something up on my behalf, because you would.”

“I would,” Sam had admitted.

“Don’t.” And then, after a moment: “Tell her I asked, though. If she comes. That’s true and it’s very small and it’s all I’ve got.”

Then he had closed his eyes and gone, and Sam had sat down on the floor beside the chair with the wax tablet on his knees and begun the longest hour and forty minutes of his life.

\sceneBreak

It went badly almost at once, and the terrible thing was that Sam knew enough to know how badly and nowhere near enough to do anything about it.

That was the Citadel’s real gift, in the end. Not healing. Sam had never once in his life healed anybody; he had read about healing extensively, in seven languages, with excellent recall. What three years at the Citadel had actually given him was the ability to sit beside a person and correctly identify, in proper terminology, the order in which they were coming apart.

\textit{Respiration shallow and irregular at forty-one to the minute.} He wrote it down because writing things down was the only action available. \textit{Pulse at the throat very fast and thready and cannot be counted reliably. Skin cold at the extremities, warm at the trunk. Eyes open, unmoving, pupils very wide and not responding to a candle.}

At about the twentieth minute the seizures started.

He had read about those too. He knew, in an exact and useless way, that what he was looking at was not painful — that the body doing that had no one in it to experience it — and he found that this piece of information was worth absolutely nothing at all when the boy you had been wheeling around a camp for two months arched off a chair in front of you with his heels drumming on the floor.

He got Bran onto the floorboards so he could not fall, and got a rolled cloak under his head, and put his own hands on his shoulders and held him there and talked to him.

He could not afterward remember one word of what he said. He was fairly sure a good deal of it had been about Gilly.

\sceneBreak

The city came up through the floor the whole time.

That was the thing Sam would never manage to convey to anybody afterward, in any account: that he spent the most important hour and forty minutes of the war in a stone room at the top of a tower, holding a boy’s shoulders, entirely blind, with the whole thing happening \textit{underneath} him as noise.

He learned to read some of it. There was a low continuous grinding that was the fighting in the lower city, and there was the crack of scorpion housings, which came in ones and twos, and around the thirtieth minute there was a sound he had never heard in his life and identified immediately, because there was only one thing it could be — an enormous impact somewhere southeast, close enough to shake grit off the ceiling and into his hair, followed a second later by the entire room going a colour it had no business going.

Orange. Real orange, coming in through the arrow-slit and lying across the floorboards and Bran’s face, wavering.

\textit{That’s dragonfire,} Sam thought. \textit{That’s actual dragonfire and it’s not moving.}

And Bran’s body, on the floor under his hands, went completely still.

\sceneBreak

The stillness was worse than the seizing.

Sam got his fingers to the boy’s throat and could not find anything, and made himself count to ten, and found it — very slow now, very far apart, a beat and then a long nothing and then another beat. The breathing had gone from forty-one to the minute to something Sam did not write down because he could not make himself write it down.

“No,” he said. “No, no, no. Bran. Bran, listen — it worked, whatever you did it \textit{worked}, I saw the light, the whole room went— you can come back now, there’s nothing left out there that needs you, come back.”

He was aware, saying it, that he had no idea whether any part of that was true.

That was the shape of the next forty minutes and it was the purest helplessness Sam Tarly ever experienced. He did not know if the plan had worked. He did not know if the Night King was dead or walking into the Mud Gate. He did not know if Jon was alive or if Arya was or if there was going to be a city in the morning. He had one piece of information — that something enormous had burned southeast of him at the thirtieth minute — and one boy on the floorboards who had stopped seizing and had not come back, and no way on earth to connect the two.

So he did what he had told Bran he would do. He sat on the floor of a tower with the tablet on his knees and he waited to write things down.

\sceneBreak

Bran spoke four times.

Sam got all of it. He would spend a considerable part of the rest of his life being quietly, fiercely proud of that — that his hand had been steady enough and his training good enough that he got every word of it exactly, with the pauses marked, on a wax tablet, in a room where he thought the world was ending.

The first came at what he estimated afterward to be an hour and five minutes.

“— carved all of them. Every tree. Every single one, I didn’t know that until I was in it —“

The voice was wrong. That was the first thing and Sam nearly dropped the stylus. It was Bran’s voice and it was coming out of Bran’s mouth and there was nobody behind it; it had the flat unweighted quality of a man reading aloud from a page he is not attending to. And the eyes did not move.

Sam wrote it down. \textit{Carved all of them. Every tree. Every single one, I didn’t know that until I was in it.}

The second came perhaps four minutes later.

“— they knew what they’d already made. They knew and they signed it anyway. That’s the part nobody —“

And stopped.

Sam wrote that down too, and then sat looking at it, because he was a scholar and had been since he was nine years old and there was a piece of him that could not be switched off by a war, and that piece of him had just noticed something and did not like it.

\textit{Signed.} Signed what? Who had signed anything? He turned it over and got nothing, and made a small mark in the wax beside it to indicate a thing he intended to think about later, which was a habit he had learned at the Citadel and had never broken.

The third came almost immediately after.

“— it isn’t a wall yet. There’s only the agreement and there’s nobody to —“

Sam’s hand stopped moving.

He sat on the floorboards of the Tower of the Hand with the stylus in the wax and read what he had just written and read it again, and the thing that had gone cold in his chest was not fear exactly, or grief, but the specific vertigo of a man who has found a mistake in a text he trusts.

\textit{Isn’t.}

Not \textit{wasn’t}. Not \textit{hadn’t been}. Sam Tarly had spent three years being beaten about the head with tenses by an old man in Oldtown who believed, correctly, that most historical error came from people being careless about them, and he knew — he \textit{knew}, with the whole of his training, the way you know your own name — that a man describing a thing that had happened in the past said \textit{wasn’t}.

\textit{Isn’t a wall yet} was what a person said standing in front of the place where one was going to be.

He looked at Bran’s face. It had not changed. The eyes were open and had not moved in an hour and there was nothing behind them at all.

“Bran,” Sam said, very quietly. “Bran, where are you?”

The fourth came about a minute after that, and it was the only one that had anything in it that sounded like a person.

“— someone should tell her I could still hear her breathing —“

Then nothing else, ever again.

\sceneBreak

Sam wrote it down. He wrote all four of them out again, properly, in ink, on paper, at some point in the following hours when his hands would do it, and he checked the ink copy against the wax three separate times, and then he did not scrape the wax, because it seemed to him that a wax tablet was a thing that could be checked against later and that people two hundred years from now might want to.

He did not understand any of it.

That was the honest situation and he was scrupulous about it for the rest of his life, in the four separate accounts he eventually wrote and in every conversation he ever had on the subject. He did not understand it at the time. He had a boy on a floor who had stopped being in his body, and four sentences that did not connect to one another or to anything else, and one wrong tense that bothered him in a way he could not explain to anybody, including himself.

\textit{Her breathing.} He had turned that one over more than the others, at the time, because it was the only one with a person in it — and because Bran had said a woman’s name to him ninety minutes earlier and then failed to finish the sentence, and Sam was not a clever man in the ways that got you noticed but he was not a stupid one either.

He had met Meera Reed exactly once, in a ruin north of the Wall, in the dark, and had given her a bundle of dragonglass and a great deal of unsolicited advice, and had spent perhaps two hours in her company. What he remembered most clearly was that she had been hauling a boy on a sledge and had not once, in two hours, complained about it or mentioned it or asked for help.

He wrote \textit{her breathing} out in ink and put a mark beside it and did not know what else to do.

\sceneBreak

It was almost light when somebody finally came.

Sansa. Not a runner — Sansa Stark herself, in the doorway, with ash to her elbows and a cut over one eyebrow that nobody had cleaned, and Sam looked up at her from the floor and knew before she said anything that she had come up eleven flights of stairs in the middle of a catastrophe because there was one thing left in this city she needed to see with her own eyes.

He did not have the words ready. He had had an hour and forty minutes to prepare them and he had spent all of it writing.

“He’s breathing,” he said. “Sansa, he’s — he’s alive, he’s breathing on his own, it’s slow but it’s steady and it’s been steady an hour. But he’s — “

He stopped.

“He’s not there,” Sansa said.

“No.”

She came into the room and got down on the floor beside her brother, in her furs, in the ash, without any of the care a great lady is supposed to take about such things, and she looked at Bran’s face for a long moment.

“Is it over?” Sam asked. “Out there. Is it — nobody’s told me anything, I’ve been up here the whole — “

“They think so.” Sansa did not look up. “Nobody’s certain. It stopped, mostly, about two hours ago, in pieces, and there’s still fighting in the north of the city because there are still things moving and nobody can tell whether they’re moving because something’s telling them to or because they haven’t finished falling over.” Her voice was absolutely level and Sam understood that it was costing her everything she had. “Jon’s alive. His hand is — they don’t think he’ll keep it. Arya’s alive; her arm’s going bad and she won’t let anyone look at it. Daenerys is alive. Ser Jaime is dead. Lord Varys is dead. Ser Roland’s lost about half his company. They’re still bringing people in and nobody has any numbers yet and I have to go back down in a moment and pretend I know what to do about all of it.”

“Sansa — “

“So tell me quickly,” she said, “because I’ve got about a hundred heartbeats before I have to be the person who’s holding this city together again. Did he say anything?”

Sam looked down at the tablet in his hands.

He thought about \textit{it isn’t a wall yet}, and the wrong tense, and the mark he had scratched in the wax beside it.

“Four things,” he said. “None of them make any sense.”

“Say them anyway.”

So he did — all four, exactly, with the pauses, the way he had written them — and Sansa Stark knelt on the floorboards beside her brother’s body and listened to every word, and at the end of it she was quiet for a moment.

“No,” she said. “None of that means anything to me either.”

“I’m sorry.”

“Don’t.” She got up, and straightened her furs, and became the other thing again in front of him — Sam watched it happen, the way he had watched it happen at four different war councils, and it was no less unsettling the fourth time. “Keep it. All of it. Keep the tablet and the copy and don’t let anybody throw either of them away, and if you write more of it down, keep that too.” She was already at the door. “You’ve spent this whole war being the only person who kept things that didn’t matter yet, Sam. Go on doing it. I’ve stopped believing anything in this family is finished simply because it’s stopped happening.”

Then she went back down the eleven flights, and Sam Tarly sat on the floor of a tower at dawn beside a boy who was breathing, with a wax tablet on his knees and a wrong verb on it that he would not understand for another twenty-six years.